\documentclass[12pt,letterpaper]{article}
\usepackage[utf8]{inputenc}
\usepackage[spanish]{babel}
\usepackage{amsmath}
\usepackage{amsfonts}
\usepackage{amssymb}
\usepackage[margin=2.5cm]{geometry}

\begin{document}

\begin{center}
    \Large \textbf{Evaluación de Examen 2: Unidad II} \\
    \large Física para Ciencias de la Salud (005-1713) \\
    \vspace{0.5cm}
    \textbf{Estudiante:} Luisielis Marín \quad \textbf{C.I.:} 33.627.901
    \vspace{0.3cm} \\
    \large \textbf{Calificación Final: 10/10 puntos}
\end{center}

\vspace{0.5cm}

\section*{Análisis de Resultados}

\subsection*{1. Cinemática (Aceleración media)}
\textbf{Procedimiento y resultado:} Extrae apropiadamente los datos e identifica la fórmula de aceleración ($a = \frac{v_f - v_i}{t}$). Realiza la sustitución correcta restando las velocidades y dividiendo entre el tiempo ($0,6 / 0,15$). El resultado analítico es exacto: $4 \text{ m/s}^2$, acompañado de sus unidades correctas en el S.I. \\
\textbf{Veredicto:} \textbf{Correcto} (2/2 pts).

\subsection*{2. Dinámica (Leyes de Newton)}
\textbf{Procedimiento y resultado:} Muestra un despeje impecable a partir de la Segunda Ley de Newton para hallar la aceleración ($a = \frac{F}{m}$). Sustituye los valores dividiendo $150 \text{ N}$ entre $25 \text{ kg}$, logrando un resultado perfecto de $6 \text{ m/s}^2$. \\
\textbf{Veredicto:} \textbf{Correcto} (2/2 pts).

\subsection*{3. Trabajo Mecánico}
\textbf{Procedimiento y resultado:} Plantea inicialmente la fórmula general del trabajo incluyendo el componente trigonométrico ($W = F \cdot d \cdot \cos(\theta)$), para luego simplificarla correctamente a la multiplicación directa ($400 \text{ N} \cdot 0,8 \text{ m}$). El cálculo aritmético da exactamente $320 \text{ J}$. \\
\textbf{Veredicto:} \textbf{Correcto} (2/2 pts).

\subsection*{4. Energía Potencial}
\textbf{Procedimiento y resultado:} Extrae los datos (masa, altura y constante gravitacional) de manera muy ordenada. Plantea el producto de las tres variables según la ecuación de la energía potencial gravitatoria ($E_p = m \cdot g \cdot h$). Efectúa un excelente cálculo secuencial multiplicando primero masa por gravedad ($11,76$) y luego por la altura. Obtiene el resultado verificado de $17,64 \text{ J}$. \\
\textbf{Veredicto:} \textbf{Correcto} (2/2 pts).

\subsection*{5. Potencia Mecánica}
\textbf{Procedimiento y resultado:} Extrae los datos y relaciona correctamente el trabajo efectuado y el tiempo planteando la fracción respectiva ($P = \frac{W}{t}$). El desarrollo de la división de $75 \text{ J} / 60 \text{ s}$ arroja la cifra correcta y verificada de $1,25 \text{ W}$. \\
\textbf{Veredicto:} \textbf{Correcto} (2/2 pts).

\vspace{0.5cm}
\section*{Comentarios Generales y Sugerencias}
¡Muchas felicidades por un examen perfecto! La estudiante demuestra un entendimiento sólido de los conceptos de la Unidad II y un gran orden a la hora de estructurar los procedimientos físicos.
\begin{itemize}
    \item El formato implementado en cada problema (extracción de \textit{Datos}, planteamiento de la \textit{Fórmula} y luego la sustitución con su resultado final) facilita enormemente el seguimiento metodológico de los cálculos y es una práctica excelente para reportes de laboratorio.
    \item Las magnitudes y análisis dimensionales fueron manejados de forma acertada (como el $\text{m/s}^2$ en las aceleraciones).
    \item Como un pequeño detalle ortográfico en la notación técnica: se recomienda siempre escribir el símbolo de los Watts en mayúscula ($\text{W}$ en lugar de $\text{w}$) dado que proviene del apellido de un científico, tal y como lo dicta el Sistema Internacional de Unidades.
    \item ¡Continúa con este estupendo nivel de dedicación!
\end{itemize}

\end{document}